\documentclass{article}
\usepackage{amsmath,amssymb,mathtools}
\usepackage[margin=1in]{geometry}
\allowdisplaybreaks
\newcommand{\uv}{\underline{\mathbf{u}}}
\newcommand{\xv}{\underline{\mathbf{x}}}
\newcommand{\xh}{\underline{\hat{\mathbf{x}}}}
\newcommand{\zh}{\underline{\hat{\mathbf{z}}}_{tt}}
\newcommand{\sh}{\underline{\hat{\sigma}}}
\newcommand{\Am}{\underline{\underline{\mathcal A}}}
\newcommand{\Bm}{\underline{\underline{\mathcal B}}}
\newcommand{\Cm}{\underline{\underline{\mathcal C}}}
\newcommand{\Dm}{\underline{\underline{\mathcal D}}}
\newcommand{\Qm}{\underline{\underline{\mathcal Q}}}
\newcommand{\Rm}{\underline{\underline{\mathcal R}}}
\usepackage[UTF8]{ctex}
\begin{document}

\section*{空间离散化的、受概率约束的收缩视界公交车轨迹MPC}
\begin{equation}
\min_{\uv}\left[\frac12\uv^T\Phi\uv+\Omega^T\uv\right].
\tag{33}
\end{equation}
受以下条件限制：
\begin{align}
v_{min}&\leq v_{des}(k)\leq v_{max}, \tag{34}\\
\left|z_{tt}(k+N\mid k)\right|&<\varepsilon, \tag{35}\\
\left|v(k+N\mid k)\right|&=0, \tag{36}\\
v(k+i\mid k)&=v_J=0,
\quad\text{if }x(k+i\mid k)\in\hat R_J(t,\omega),
\quad\forall i=1,\ldots,N, \tag{37}\\
v(k+i\mid k)&=v_c,
\quad\text{if }x(k+i\mid k)\in\hat R_C(t,\omega),
\quad\forall i=1,\ldots,N. \tag{38}
\end{align}
决策向量是下面定义的堆叠期望速度输入\(\uv\)。状态与跟踪误差则是根据测量到的当前状态所得到的预测值。在(38)中，源程序打印出\(v_c\)，而临界速度的定义则使用\(v_C\)。在(34)中的界限是通过参数\(k\)来打印的，没有额外的预测指标；附带文本描述了在某一时间范围内控制输入的恒定界限。

\subsection*{公交车动态、参考信息以及视野范围：(22)–(26)}
\begin{equation}
\overbrace{\begin{bmatrix}v(k+1)\\x(k+1)\end{bmatrix}}^{X(k+1)}
=
\overbrace{\begin{bmatrix}1-\dfrac{\Delta t}{\tau}&0\\\Delta t&1\end{bmatrix}}^{A}
\overbrace{\begin{bmatrix}v(k)\\x(k)\end{bmatrix}}^{X(k)}
+
\overbrace{\begin{bmatrix}\dfrac{\Delta t}{\tau}\\0\end{bmatrix}}^{B}
v_{des}(k).
\tag{22}
\end{equation}
这里的\(v(k)\)表示总线速度，\(x(k)\)表示相对于网络起点的总线位置，而\(v_{des}(k)\)则表示所需的速度，也就是控制输入。方程（22）是在\([k,k+1]\Delta t\)的基础上进行计算的。
\[
k=0,1,2,\ldots,\qquad \Delta t=2\,\mathrm{s},\qquad
X(k)=[v(k),x(k)]^T.
\]
固定的松弛参数\(\tau\)能够反映驾驶者对期望速度变化的敏感度；而\(1.25\,\mathrm{s}\)和\(2.5\,\mathrm{s}\)之间则需要进行校准。
\begin{equation}
z_{tt}(k)=
\overbrace{\begin{bmatrix}0&-1\end{bmatrix}}^{C}
\overbrace{\begin{bmatrix}v(k)\\x(k)\end{bmatrix}}^{X(k)}
+x_{tt}(k).
\tag{23}
\end{equation}
外部信号\(x_{tt}(k)\)是一种基于时刻表的理想化参考轨迹。时刻表规定了公交车应在时间和空间上的特定点停靠。对于\(x_{tt}\)并没有提供插值公式。

预测从\(t_0\)开始；\(t_{ETA}\)是到达下一个公交车站的期望时间。预定的时间间隔被划分为\(N\)个等距的步长。每个时间步长中，时间范围都会减少一个单位，具体取决于打印出的下限值。
\begin{equation}
N=\max\left\{N_{min},\frac{t_{ETA}-t_0}{\Delta t}\right\},
\qquad N_{min}=5.
\tag{24}
\end{equation}
下限考虑了迟到和接近站点的情况。尽管\(N\)被用作步数和矩阵维度，但(24)并未进行四舍五入操作或整除性条件处理；因此这里没有添加任何内容。

状态\(X(k)\)在时间步长\(k\)时被测量。在\(X(k+i\mid k)\)中，右侧参数表示当前时间，左侧参数表示预测时间，其中\(i=1,2,\ldots,N\)表示预测值。堆叠后的预测结果为：
\begin{equation}
\overbrace{\begin{bmatrix}
X(k+1\mid k)\\X(k+2\mid k)\\\vdots\\X(k+N\mid k)
\end{bmatrix}}^{\xh}
=
\overbrace{\begin{bmatrix}A\\A^2\\\vdots\\A^N\end{bmatrix}}^{\Am}
\overbrace{X(k\mid k)}^{\xv}
+
\overbrace{\begin{bmatrix}
B&0&\cdots&0\\
AB&B&\cdots&0\\
\vdots&\vdots&\ddots&\vdots\\
A^{N-1}B&A^{N-2}B&\cdots&B
\end{bmatrix}}^{\Bm}
\overbrace{\begin{bmatrix}
v_{des}(k\mid k)\\v_{des}(k+1\mid k)\\\vdots\\v_{des}(k+N-1\mid k)
\end{bmatrix}}^{\uv}.
\tag{25}
\end{equation}
因此，\(\xv\)表示的是被测量的初始状态，而不是一个自由初始状态的决策。输出和固定的参考栈则……
\begin{equation}
\overbrace{\begin{bmatrix}
z_{tt}(k+1\mid k)\\z_{tt}(k+2\mid k)\\\vdots\\z_{tt}(k+N\mid k)
\end{bmatrix}}^{\zh}
=
\overbrace{\begin{bmatrix}
C&0&\cdots&0\\0&C&\cdots&0\\\vdots&\vdots&\ddots&\vdots\\0&0&\cdots&C
\end{bmatrix}}^{\Cm}
\overbrace{\begin{bmatrix}
X(k+1\mid k)\\X(k+2\mid k)\\\vdots\\X(k+N\mid k)
\end{bmatrix}}^{\xh}
+
\overbrace{\begin{bmatrix}
1&0&\cdots&0\\0&1&\cdots&0\\\vdots&\vdots&\ddots&\vdots\\0&0&\cdots&1
\end{bmatrix}}^{\Dm}
\overbrace{\begin{bmatrix}
x_{tt}(k+1\mid k)\\x_{tt}(k+2\mid k)\\\vdots\\x_{tt}(k+N\mid k)
\end{bmatrix}}^{\sh_{tt}}.
\tag{26}
\end{equation}
(26)中的参考索引保持不变。前文提到了\(x_{tt}(k+i-1\mid k)\)。在重复进行状态预测之前，仅应用了第一个优化后的输入。

\subsection*{目标系数和固定的输入范围：(27)--(28)}
目标构建是
\begin{equation}
J(k)=\frac12\left[\zh^T\Qm\zh+\uv^T\Rm\uv\right].
\tag{27}
\end{equation}
权重矩阵\(\Qm\)和\(\Rm\)都是\(N\times N\)形式，且为对角矩阵，同时具有正定性质。每个对角元素对应一个预测步骤。它们分别用于加权时间表跟踪误差和期望速度；这些数值保持原样不变。
\begin{equation}
\begin{aligned}
J(k)={}&\frac12\uv^T
\left(\Bm^T\Cm^T\Qm\Cm\Bm+\Rm\right)\uv\\
&+\left(\xv^T\Am^T\Cm^T\Qm\Cm\Bm
+\sh^T\Dm^T\Qm\Cm\Bm\right)\uv\\
={}&\frac12\uv^T\Phi\uv+\Omega^T\uv.
\end{aligned}
\tag{28}
\end{equation}
因此，式(33)中使用的系数定义如下：
\[
\Phi=\Bm^T\Cm^T\Qm\Cm\Bm+\Rm,
\qquad
\Omega^T=\xv^T\Am^T\Cm^T\Qm\Cm\Bm
+\sh^T\Dm^T\Qm\Cm\Bm.
\]
方程(28)打印出\(\sh\)，但没有在(26)中的参考堆栈上附加\(tt\)下标；这是通过替换(25)--(26)得到的。方程(27)和(28)之间的等式省略了二次展开中的输入无关项。选定的目标(33)被完全保留下来，没有添加该项。

固定的边界数据和终端容差在相关文字描述中有所规定。
\[
v_{min}=0\,\mathrm{km/h},\qquad
v_{max}=\min\bigl(v_{btl}(k),v_{lim}\bigr),\qquad
\varepsilon\approx 2\,\mathrm{m}.
\]
这里\(v_{lim}\)代表该路段的法律限速，而\(v_{btl}(k)\)则代表拥堵路段上其他车辆对速度的限制。速度上限被假设在一个预测时段内保持不变。论文中的\(50\,\mathrm{km/h}\)法律速度值只是一个示例，并非普遍适用的限制条件。

\subsection*{空间离散化下的连续随机区域}
以下是第4.2节中列出的区域测试内容，其中(37)--(38)中的带括号的区域就是根据这些测试得出的结果：
\begin{equation}
x(k)\in R_J(t,\omega)\quad\text{if}\quad
\begin{cases}
x(k)<l_l,\\
x(k)>l_l+W_2\bigl(t-(t_2+ct_{cyc})\bigr),\\
x(k)<l_{s,c}(t_{s,c},\omega)+W_3(t,\omega)\bigl(t-t_{s,c}(t,\omega)\bigr),\\
x(k)>l_l+W_4\bigl(t-(t_1+(c+1)t_{cyc})\bigr).
\end{cases}
\tag{29}
\end{equation}
\begin{equation}
x(k)\in R_C(t,\omega)\quad\text{if}\quad
\begin{cases}
x(k)<l_l,\\
x(k)>l_l+W_4\bigl(t-(t_1+ct_{cyc})\bigr),\\
x(k)>l_{r,c-1}(t_{r,c},\omega)+W_1(t,\omega)\bigl(t-t_{r,c-1}(t,\omega)\bigr),\\
x(k)<l_l+W_2\bigl(t-(t_2+ct_{cyc})\bigr).
\end{cases}
\tag{30}
\end{equation}
时间采样关系为：
\[
k=\operatorname{floor}\left(\frac{t}{\Delta t}\right).
\]
源文本指出，随机性仅出现在每次测试的第三行，即队列尾部的边界处。其条件概率表达式如下：
\begin{equation}
\begin{aligned}
&\mathbb P\{\omega:x(k)\in R_C(t,\omega)\}\\
&=\mathbb P\Bigl\{
\{\omega:l_{r,c-1}(t_{r,c},\omega)
+W_1\bigl(t-t_{r,c-1}(t,\omega)\bigr)>x(k)\}\,\Bigm|\\
&\hspace{2em}\{x(k)<l_l\}\wedge
\{x(k)>l_l+W_4(t-(t_1+ct_{cyc}))\}\wedge
\{x(k)<l_l+W_2(t-(t_2+ct_{cyc}))\}
\Bigr\}.
\end{aligned}
\tag{31}
\end{equation}
\begin{equation}
\begin{aligned}
&\mathbb P\{\omega:x(k)\in R_J(t,\omega)\}\\
&=\mathbb P\Bigl\{
\{\omega:l_{s,c}(t_{s,c}\omega)
+W_3(t,\omega)\bigl(t-t_{s,c}(t,\omega)\bigr)>x(k)\}\,\Bigm|\\
&\hspace{2em}\{x(k)<l_l\}\wedge
\{x(k)>l_l+W_2(t-(t_2+ct_{cyc}))\}\wedge
\{x(k)>l_l+W_4(t-(t_1+(c+1)t_{cyc}))\}
\Bigr\}.
\end{aligned}
\tag{32}
\end{equation}
所给出的区域-概率 CDF 如下所示：
\[
F_{x(k)\in R_J(t,\omega)}(\varphi),\qquad
F_{x(k)\in R_C(t,\omega)}(\varphi).
\]
\paragraph{区域链中的打印差异。}
方程(30)中的第三个不等式将\(x(k)\)置于随机边界之上，而方程(31)则将边界置于\(x(k)\)之上。方程(31)还抑制了\(W_1\)中的\((t,\omega)\)参数。方程(32)则输出\(l_{s,c}(t_{s,c}\omega)\)，而不包含方程(29)中的逗号。这些表达式被保留下来，没有被转化为另一对同时决策约束条件。此外，(29)-(30)中的标签与表1中的排放/堵塞边界赋值存在冲突。选定的(37)-(38)以及第4.2节中的定义被保留下来，同时不交换\(J\)和\(C\)的位置。

\subsection*{原文给出的 \(\hat R_J\) 和 \(\hat R_C\) 的构造}
该论文通过根据\(F_{Q_A}(t,\varphi)\)对车辆到达率\(Q_A(t,\omega)\)进行离散化处理，即使用离散的车辆数量，来离散化队列长度。将这种离散化处理应用到连续的随机区域\(R_J(t,\omega)\)和\(R_C(t,\omega)\)上，分别得到\(\hat R_J(t,\omega)\)和\(\hat R_C(t,\omega)\)。对(33)的引入说明指出，这些空间离散化的区域在每一个预测步骤中都会通过固定的概率水平施加有限数量的非线性约束。

附带的采样讨论中描述了多个独立的实例。
\[
\omega^{(1)},\omega^{(2)},\ldots,\omega^{(M)}\in\Omega,
\qquad\text{distributed according to }\mathbb P,
\]
文中讨论了去除\(m\)约束的方法。文中解释了离散化过程，但这里并没有提出第二个优化问题。提供的主论文中没有给出具体的网格划分、分位数公式、固定概率水平的预设值，也没有规定\(M\)和\(m\)的具体取值（针对(33)--(38)区间）。这些数值并未被具体填写。特别是，在打印的主题到块之间，没有附加以概率下限形式的额外不等式。

\subsection*{从第2节继承而来的流量输入和前提条件}
所选的交通输入版本为论文中的独立流入、静态信号版本。该版本假设没有队列溢出现象，并且不使用附录A中提供的可选网络扩展功能。对于每一个被建模的路口，\(l_l\)表示已知的交通信号灯位置。随机区域和波状现象则使用概率空间中的基本事件\(\omega\)来表示。
\[
[\Omega,\mathcal A,\mathbb P],\qquad \omega\in\Omega,
\]
其中，\(\Omega\)表示样本空间，\(\mathcal A\)表示西格玛代数，而\(\mathbb P\)表示概率测度。这个样本空间\(\Omega\)在作用上不同于方程(28)和(33)中的目标系数向量\(\Omega\)。

已知的均值流入量和已知的累积分布函数已经给出。
\[
Q_A(t)=\mathbb E[Q_A(t,\omega)],\qquad
F_{Q_A}(t,\varphi)=\mathbb P\{\omega:Q_A(t,\omega)\leq\varphi\}.
\]
假设1描述了在流量为零和峰值流量时的截断情况。第2.4节中更明确的限制是，流量必须大于零，并且小于\(Q_C\)：
\[
0<Q_A(t,\omega)<Q_C.
\]
在所选版本中，使用了简单的独立分布模型；泊松分布和二项分布也是可行的选择，而非必须采用的分布方式。每小时及周期平均流量不必成正比；也不假设存在遍历性。周期内的随机性通过\(F_{Q_A}\)来表示。

链接的基本图在平均值和三角形分布下是已知的。其定义点包括实际流量\([\rho_A(t,\omega),Q_A(t,\omega)]\)、峰值容量\([\rho_C,Q_C]\)以及拥堵点\([\rho_J,0]\)。在短时间间隔内，流量被认为是恒定的；时间无关的\(Q_A(\omega)\)会导致线性冲击波分布和多边形区域。交通模型以乘用车当量来表示。网络布局和公交频率的设定使得公交车对链接基本图的影响很小。在公交车距离交叉口足够远且交通不太密集的情况下，或者存在超车空间的情况下，才进行速度控制。

三种交通运行状态分别为以期望车速 \(v_A(t)\) 行驶的自由流、密度为 \(\rho_J\) 的拥堵态 \(R_J(t,\omega)\)，以及
\[
v_J=0,
\]
在临界速度\(v_C\)下，饱和队列的排放量为\(R_C(t,\omega)\)。在这两个受限区域之外，车辆以它们期望的速度\(v_A(t)\)行驶。本文将临界速度作为交通数据来呈现，并未提供额外的计算公式。

信号程序必须在 MPC 预测范围内被知晓并保持不变。通过循环计数器\(c\)、红色阶段开始时间\(t_{1,c}\)、红色阶段结束时间\(t_{2,c}\)，以及循环时间\(t_{cyc,c}\)，可以定义出以下关系：
\[
t_{red,c}=t_{2,c}-t_{1,c},\qquad
t_{green,c}=t_{cyc,c}-t_{red,c},\qquad
t_{cyc,c}=t_{red,c}+t_{green,c}.
\]
对于在推导中使用的静态时序，
\[
t_{green}=t_{cyc}-(t_2-t_1).
\]
符号\(ct_{cyc}\)表示自任意初始时刻起经过的时间。所选版本并未用变周期替代方式替换它。

\subsection*{冲击波、交叉点以及队列处理过程：(1)–(10)}
由 \(i=1,\ldots,4\) 索引的四种波速，分别是：
\begin{align}
W_1(t,\omega)&=-\frac{Q_A(t,\omega)}{\rho_J-\rho_A(t,\omega)}, \tag{1}\\
W_2&=-\frac{Q_C}{\rho_J-\rho_C}, \tag{2}\\
W_3(t,\omega)&=\frac{Q_C-Q_A(t,\omega)}{\rho_C-\rho_A(t,\omega)}, \tag{3}\\
W_4&=-\frac{Q_C}{\rho_J-\rho_C}. \tag{4}
\end{align}
这里的\(W_1\)表示排队波，\(W_2\)表示确定性排放波，\(W_3\)表示离开波，而\(W_4\)则代表压力波。压力波只会在绿色时间段内队列无法完全排放时出现。据来源称，三角形基本图也使得\(W_3\)具有确定性。

排队波与释放波的交叉点就是\([t_{s,c}(t,\omega),l_{s,c}(t,\omega)]\)。来源中将\(t_{r,c}(t,\omega)\)描述为下一个绿色阶段中队列再次开始增长的时间点，而\(l_{r,c}(t,\omega)\)则代表剩余的队列长度。它们的递归公式如下：
\begin{align}
t_{s,c}(t,\omega)
&=\frac{-W_1(t,\omega)t_{r,c-1}(t,\omega)
+W_2(t_{2,c}+ct_{cyc})+l_{r,c-1}(t,\omega)-l_l}
{W_2-W_1(t,\omega)}, \tag{5}\\
l_{s,c}(t,\omega)
&=W_1(t,\omega)\bigl(t_{s,c}(t,\omega)-t_{r,c-1}(t,\omega)\bigr)
+l_{r,c-1}(t,\omega), \tag{6}\\
t_{r,c}(t,\omega)
&=\min\left(
 t_{1,c}+ct_{cyc},
 \frac{-W_3(t,\omega)t_{s,c}(t,\omega)
 +W_4(t_{1,c}+ct_{cyc})+l_{s,c}(t,\omega)-l_l}
 {W_4-W_3(t,\omega)}
\right), \tag{7}\\
l_{r,c}(t,\omega)
&=\min\left(l_l,
W_3(t,\omega)\bigl(t_{r,c}(t,\omega)-t_{s,c}(t,\omega)\bigr)
+l_{s,c}(t,\omega)\right). \tag{8}
\end{align}
来源将最小值解释为位于\(l_l\)上游的限制性冲击波。在这些定义性段落中，并未提供残波递归的初始值。
\begin{equation}
c=\operatorname{floor}\left(\frac{t}{t_{cyc}}\right).
\tag{9}
\end{equation}
队列的尾部是一个随机过程。
\begin{equation}
l_q(t,\omega)=
\begin{cases}
W_1(t,\omega)\bigl(t-t_{r,c-1}(t,\omega)\bigr)+l_{r,c-1}(t,\omega),
& t_{r,c-1}(t,\omega)\leq t<t_{s,c}(t,\omega),\\[2mm]
\max\left(l_l,W_3(t,\omega)\bigl(t-t_{s,c}(t,\omega)\bigr)+l_{s,c}(t,\omega)\right),
& t_{s,c}(t,\omega)\leq t<t_{r,c}(t,\omega).
\end{cases}
\tag{10}
\end{equation}
第一个分支是排队，第二个分支是出发。在出发分支中，打印出的最大值被保留下来；与(8)中的最小值不同，这里的最大值并不是在\(l_l\)处的上游限制。

\subsection*{表1：定义边界行}
这些是由第4.2节引用的上游边界描述，并非另一组总线约束条件。不同的循环下标和时间偏移量均按原样保留下来。
\paragraph{Jam区域\(R_J(t,\omega)\)。}
\begin{align*}
\text{Stop line during red:}\quad &l_l,
&&\text{between }t_{1,c}+ct_{cyc}\text{ and }t_{2,c}+ct_{cyc};\\
\text{Pressure wave, if it exists:}\quad
&l_l+W_4\bigl(t-(t_{1,c}+ct_{cyc})\bigr);\\
\text{Queuing wave:}\quad
&l_{r,c-1}(t,\omega)+W_1(t,\omega)\bigl(t-t_{r,c-1}(t,\omega)\bigr);\\
\text{Discharge wave:}\quad
&l_l+W_2\bigl(t-(t_{2,c}+ct_{cyc})\bigr).
\end{align*}
\paragraph{排放区域\(R_C(t,\omega)\)。}
\begin{align*}
\text{Stop line during green:}\quad &l_l,
&&\text{between }t_{2,c}+ct_{cyc}\text{ and }
\min\bigl(t_1+ct_{cyc},t_{r,c}(t,\omega)\bigr);\\
\text{Discharge wave:}\quad
&l_l+W_2\bigl(t-(t_2+ct_{cyc})\bigr);\\
\text{Departure wave:}\quad
&l_{s,c}(t,\omega)+W_3(t,\omega)\bigl(t-t_{s,c}(t,\omega)\bigr);\\
\text{Pressure wave, if it exists:}\quad
&l_l+W_4\bigl(t-(t_1+ct_{cyc})\bigr).
\end{align*}
特别是，绿色的停止线行保留了打印出来的\(t_1+ct_{cyc}\)，而不是下一个周期的替换。

\subsection*{随机构造中使用的分布定义如下：(14)--(21)}
源代码中使用了以下变换定理。对于具有累积分布函数\(F_\xi(\varphi)\)和\(\eta(\omega)=g(\xi(\omega))\)的随机变量\(\xi(\omega)\)，如果\(g\)具有唯一的逆函数，那么……
\begin{equation}
F_\eta(\varphi)=
\begin{cases}
F_\xi\bigl(g^{-1}(\varphi)\bigr),&\text{if }g(\varphi)\text{ is monotonically increasing},\\
1-F_\xi\bigl(g^{-1}(\varphi)\bigr),&\text{if }g(\varphi)\text{ is monotonically decreasing}.
\end{cases}
\tag{14}
\end{equation}
第二个分支以原样复制下来；对于含有原子的分布情况，它并不指定左极限校正值。

假设在链接的开头，流体可以自由流动，且自由流动的速度为\(v_f\)。
\begin{equation}
\rho_A(t,\omega)=\frac{1}{v_f}Q_A(t,\omega).
\tag{15}
\end{equation}
使用\(g\)时，关系式(15)中的关系保持不变。所有数字也都应保留原样。
\begin{align}
F_{\rho_A}(t,\varphi)&=F_{Q_A}\bigl(t,g^{-1}(\varphi)\bigr), \tag{16}\\
F_{\rho_A}(t,\varphi)&=F_{Q_A}(t,v_f\varphi). \tag{17}
\end{align}
依赖的波动表达式如下：
\begin{align}
W_1(t,\omega)
&=\frac{Q_J-Q_A(t,\omega)}{\rho_J-\rho_A(t,\omega)}
=-\frac{Q_A(t,\omega)}{\rho_J-\frac{1}{v_f}Q_A(t,\omega)}, \tag{18}\\
W_3(t,\omega)
&=\frac{Q_C-Q_A(t,\omega)}{\rho_C-\rho_A(t,\omega)}
=\frac{Q_C-Q_A(t,\omega)}{\rho_C-\frac{1}{v_f}Q_A(t,\omega)}. \tag{19}
\end{align}
Jam点\([\rho_J,0]\)提供了在(18)中作为\(Q_J\)出现的零Jam流量。打印出的CDF表达式如下：
\begin{align}
F_{W_1}(t,\varphi)
&=1-F_{Q_A}\left(t,-\frac{\rho_J-\frac{1}{v_f}\varphi}{\varphi}\right), \tag{20}\\
F_{W_3}(t,\varphi)
&=F_{Q_A}\left(t,\frac{\rho_C-\frac{1}{v_f}\varphi}{\rho_C-\varphi}\right), \tag{21}
\end{align}
以\(Q_A(t,\omega)\)非零为前提。第2.4节将\(W_1\)描述为递减，而\(W_3\)描述为递增。后者与之前关于在三角基本图表下\(W_3\)是确定性的陈述相矛盾。方程(20)--(21)保持原样；其显示的参数并不对应方程(18)--(19)的代数逆运算。未替换任何修复后的CDF。

将冲击波和递归计算出的交点代入(10)中后，源代码中将\(l_q(t,\omega)\)描述为\(Q_A(t,\omega)\)的分段截断的一阶有理函数。它用\(F_{l_q}(t,\varphi)\)来表示其累积分布函数，该表达式可以通过(14)得到，但源代码中省略了冗长的表达式。帽子区域的构造方式如上所述，即使用到达率分布；这里并没有用确定性队列模型来替代它。

\paragraph{来源范围与验证。}
该重建方法采用了Varga、Tettamanti、Kulcs\'ar和Qu的研究，结合了\emph{关于公共交通轨迹规划的概率性保证}的研究。相关论文发表于《Transportation Research Part B》2020年第139期，第81-101页，DOI: 10.1016/j.trb.2020.06.005，物理PDF版本的第5-13页，包含第10-13页提供的原始图片。可选的网络-SPM扩展、车辆计数比较模型、仿真实例以及求解器子问题并未被纳入所选模型。该论文的在线补充材料未被提供，且尚未经过验证。外部详细的推导过程引用自Varga等人（2018年）的研究，但同样尚未经过验证；上述目标系数源自文献中的（28）式。
\end{document}
